\section{Conclusion}
\label{sec:conclusion}
We studied DR grading, lesion segmentation and lesion detection as one coupled problem on the DDR dataset, with the grade computed from zone-wise soft lesion statistics of a segmentation network and complemented by ordinal soft targets, a clinical-rule consistency loss and conformal referral sets. Under a strictly controlled protocol --- a leak-controlled test set, five seeds, bootstrap intervals, matched baselines and component ablations --- the clearest finding is that lesion evidence adds information that a generic image embedding lacks: the quadratic-weighted kappa rises from 0.64 to 0.73 and the AUC for referable DR from 0.87 to 0.92, with a segmentation network trained on only 383 pixel-annotated images. Soft ordinal targets improve calibration; the token architecture gives a small gain over naive fusion; the rule-consistency loss did not change accuracy or violation rates measurably, because the rules are rarely violated to begin with. Conformal prediction delivers its guarantee under exchangeability but is not robust to the shift we found between the DDR validation and test partitions, and the pipeline does not transfer zero-shot to a low-resolution external set. Segmentation of small lesions (microaneurysms) and detection at strict overlap remain weak at the resolution that our compute allowed, and mild and severe DR are poorly recognised.

These conclusions come with the caveats of a CPU-bound, frozen-backbone study with single segmentation runs. The most valuable next steps are end-to-end fine-tuning with a stronger backbone, using grade labels as weak supervision for the 95\% of images without lesion masks, anatomically defined zones and stricter rules, and prospective multi-centre validation with calibration on the deployment distribution. We release code and every intermediate result so that these steps can be taken without re-running the CPU pipeline.
